The technologies that were supposed to abolish distance---container shipping,
fibre-optic cables, the internet, cheap flights---have made the world more
integrated and faster-moving. This paper argues that for the most remote and
smallest economies these same forces turn remoteness from a one-off cost into
a \emph{recurring tax on participation}, paid on every shipment, byte,
specialist and negotiation. Using a reproducible pipeline of public data for
\nCountriesRemote{} economies, I separate being \emph{far} (GDP-weighted
distance to world output) from being \emph{small} (population). Distance, not
size, carries the income penalty: conditional on governance, doubling
remoteness is associated with \regDoubleRemotePctAbs\% lower income per head,
while population matters little. The evidence on a \emph{widening} gap is
mixed and reported as such: access to the internet has converged sharply,
and shipping gaps with the old hubs have narrowed; but connectivity at the
edge has stagnated relative to the median country, intensive use of new
technology remains far apart, and small islands lose a majority of their
graduates. Iceland, the Faroes, Singapore and Mauritius show that the edge is
not destiny---though, notably, the richest of them are small rather than
far. Remoteness is best understood as a trajectory that policy can bend:
through pooled scale, redundancy, vulnerability-based finance and circular
talent.
